% =====================================================================
%  1bat-ud3-7.tex  —  HORA 7: Funcions quadràtiques: la paràbola
%  (sense preàmbul: s'inclou des de main.tex)
% =====================================================================
\horatitol{Sessió 7 — Funcions quadràtiques: la paràbola}%
{D'on surt una corba amb un màxim: l'equació de la paràbola, el vèrtex com a punt òptim i els punts de tall.}

\subsection{Idea clau}

Hi ha situacions que ja no es modelitzen amb una recta. Pensem en el benefici d'una
venda: si apugem el preu, al principi guanyem més, però si l'apugem massa, la gent
deixa de comprar i el benefici \textbf{baixa}. La gràfica puja, fa un \textbf{cim} i
baixa: és una \textbf{paràbola}, la gràfica d'una \textbf{funció quadràtica}.

\begin{idea}
\textbf{Funció quadràtica:} $\quad f(x)=ax^2+bx+c \quad (a\neq 0)$
\begin{itemize}[leftmargin=1.4em, itemsep=2pt]
  \item La seva gràfica és una \textbf{paràbola}. Si $a>0$ s'obre cap amunt
        ($\cup$, té un \textbf{mínim}); si $a<0$ s'obre cap avall
        ($\cap$, té un \textbf{màxim}).
  \item \textbf{Vèrtex:} el punt més alt (o més baix) de la corba. La seva abscissa és
        \[
        x_V=\frac{-b}{2a},
        \]
        i la seva ordenada s'obté calculant $f(x_V)$.
  \item \textbf{Punts de tall amb l'eix $X$:} es resol $ax^2+bx+c=0$ (equació de segon
        grau de la UD2). Poden ser dos, un o cap.
  \item \textbf{Tall amb l'eix $Y$:} es fa $x=0$; és el punt $(0,c)$.
\end{itemize}
\textbf{Lectura social:} en molts problemes el vèrtex és el \textbf{punt òptim} (el
preu que dóna el màxim benefici, el nombre d'inscrits que minimitza el cost…).
\end{idea}

\subsection{Exemples}

\begin{exemple}[una paràbola amb màxim: el benefici]
El benefici d'una parada solidària segons el preu $x$ (en euros) és
$B(x)=-x^2+6x$. Com que $a=-1<0$, la paràbola s'obre cap avall i té un \textbf{màxim}.

\begin{center}
\begin{tikzpicture}
\begin{axis}[udaxis, width=9cm, height=5.6cm,
    xlabel={\small preu ($x$)}, ylabel={\small benefici},
    xmin=-0.5, xmax=6.7, ymin=-1, ymax=11,
    xtick={0,1,2,3,4,5,6}, ytick={0,3,6,9}]
  \addplot[udblue, domain=0:6, samples=80]{-x^2+6*x};
  \addplot[udnavy, only marks, mark=*, mark size=1.7pt] coordinates {(3,9) (0,0) (6,0)};
  \node[udnavy, font=\scriptsize, anchor=south] at (axis cs:3,9) {vèrtex $(3,9)$};
  \node[udnavy, font=\scriptsize, anchor=north west] at (axis cs:0,0) {$(0,0)$};
  \node[udnavy, font=\scriptsize, anchor=north east] at (axis cs:6,0) {$(6,0)$};
\end{axis}
\end{tikzpicture}

\figcap{$B(x)=-x^2+6x$: talla l'eix a $x=0$ i $x=6$; el vèrtex $(3,9)$ és el benefici màxim.}
\end{center}

\noindent El vèrtex està a $x_V=\dfrac{-6}{2\cdot(-1)}=3$, amb $B(3)=-9+18=9$. És a dir,
el preu $3\eur$ dóna el benefici màxim de $9$ (el \textbf{punt òptim}).
\end{exemple}

\begin{exemple}[una paràbola amb mínim: el cost]
El cost d'organitzar una activitat segons el nombre de monitors $x$ és
$C(x)=x^2-4x+7$. Ara $a=1>0$: la paràbola s'obre cap amunt i té un \textbf{mínim}.
El vèrtex és a $x_V=\dfrac{-(-4)}{2\cdot 1}=2$, amb $C(2)=4-8+7=3$. Amb $2$ monitors
el cost és mínim ($3$). Aquí el vèrtex és el punt de \textbf{menor cost}.
\end{exemple}

\subsection{Exercicis resolts}

\begin{exresolt}[vèrtex i talls d'una paràbola]
Per a $f(x)=-x^2+4x+5$, troba el vèrtex, els talls amb l'eix $X$ i el tall amb
l'eix $Y$. Digues si el vèrtex és màxim o mínim.
\solucio
Com que $a=-1<0$, la paràbola s'obre cap avall: el vèrtex serà un \textbf{màxim}.

\textbf{Vèrtex:}
\[
x_V=\frac{-b}{2a}=\frac{-4}{2\cdot(-1)}=2,\qquad
f(2)=-4+8+5=9.
\]
Vèrtex $(2,9)$, que és un màxim.

\textbf{Talls amb l'eix $X$} ($f(x)=0$): resolem $-x^2+4x+5=0$, o bé
$x^2-4x-5=0$. Factoritzant (o amb la fórmula): $(x-5)(x+1)=0$, d'on $x=5$ o $x=-1$.
Talls: $(5,0)$ i $(-1,0)$.

\textbf{Tall amb l'eix $Y$} ($x=0$): $f(0)=5$. Tall: $(0,5)$.
\resposta{Vèrtex $(2,9)$ (màxim); talls amb $X$: $(-1,0)$ i $(5,0)$; tall amb $Y$: $(0,5)$.}
\end{exresolt}

\begin{exresolt}[intervals de creixement i decreixement]
Per a la paràbola $f(x)=x^2-6x+8$ (que s'obre cap amunt), troba el vèrtex i digues
en quin interval creix i en quin decreix.
\solucio
$a=1>0$, s'obre cap amunt (té un mínim). Vèrtex:
\[
x_V=\frac{-(-6)}{2\cdot 1}=3,\qquad f(3)=9-18+8=-1.
\]
Vèrtex $(3,-1)$. En una paràbola cap amunt, la corba \textbf{decreix} fins al vèrtex
i \textbf{creix} a partir d'ell:
\[
\text{decreix per } x<3, \qquad \text{creix per } x>3.
\]
\resposta{Vèrtex $(3,-1)$; decreix si $x<3$ i creix si $x>3$.}
\end{exresolt}

\subsection{Exercicis per fer}

\begin{exercicis}
\begin{exlist}
  \item Per a cada paràbola, digues si s'obre cap amunt o cap avall i si el vèrtex és
        un màxim o un mínim (mira només el signe de $a$):
    \begin{enumerate}[label=\alph*), leftmargin=1.6em, topsep=2pt]
      \item $f(x)=2x^2-3x+1$
      \item $f(x)=-x^2+5$
      \item $f(x)=-3x^2+2x$
    \end{enumerate}
  \item Troba el vèrtex de cada paràbola amb la fórmula $x_V=\dfrac{-b}{2a}$:
    \begin{enumerate}[label=\alph*), leftmargin=1.6em, topsep=2pt]
      \item $f(x)=x^2-4x+3$
      \item $f(x)=-x^2+2x+8$
    \end{enumerate}
  \item Troba els punts de tall amb l'eix $X$ de $f(x)=x^2-5x+6$ (resol l'equació de
        segon grau) i el tall amb l'eix $Y$.
  \item Observa la paràbola de la gràfica. Indica, llegint-la: el vèrtex, si és màxim
        o mínim, i els punts de tall amb l'eix horitzontal.

        \begin{center}
        \begin{tikzpicture}
        \begin{axis}[udaxis, width=8.6cm, height=5cm,
            xlabel={\small $x$}, ylabel={\small $y$},
            xmin=-1.5, xmax=5.5, ymin=-5, ymax=5,
            xtick={-1,0,1,2,3,4,5}, ytick={-4,-2,0,2,4}]
          \addplot[udblue, domain=-1:5, samples=80]{x^2-4*x};
        \end{axis}
        \end{tikzpicture}

        \figcap{Llegeix el vèrtex i els talls directament de la gràfica.}
        \end{center}

  \item \textbf{(Optimització)} El benefici d'una rifa segons el preu del bitllet $x$
        és $B(x)=-2x^2+12x$. Troba el preu que dóna el \textbf{benefici màxim} (l'abscissa
        del vèrtex) i quin és aquest benefici màxim.
  \item \textbf{(Interpretació)} En una paràbola de cost $C(x)=x^2-8x+20$, què
        significa, en el context d'un problema social, el punt vèrtex? Per què sovint
        és el punt que ens interessa trobar?
\end{exlist}
\end{exercicis}

% ---------------------------------------------------------------------
\fullsolucions{Sessió 7}
\begin{solucions}
\begin{sollist}
  \item (a) Cap amunt ($a=2>0$), mínim. \quad (b) Cap avall ($a=-1<0$), màxim. \quad
        (c) Cap avall ($a=-3<0$), màxim.
  \item (a) $x_V=\dfrac{4}{2}=2$, $f(2)=4-8+3=-1$: vèrtex $(2,-1)$. \quad
        (b) $x_V=\dfrac{-2}{-2}=1$, $f(1)=-1+2+8=9$: vèrtex $(1,9)$.
  \item $x^2-5x+6=0 \Rightarrow (x-2)(x-3)=0 \Rightarrow x=2,\ x=3$. Talls amb $X$:
        $(2,0)$ i $(3,0)$. Tall amb $Y$: $(0,6)$.
  \item Vèrtex $(2,-4)$, que és un \textbf{mínim} (s'obre cap amunt); talls amb l'eix
        horitzontal: $(0,0)$ i $(4,0)$.
  \item $x_V=\dfrac{-12}{2\cdot(-2)}=3$. El preu òptim és $3\eur$, amb benefici màxim
        $B(3)=-18+36=18$.
  \item El vèrtex és el punt de \textbf{cost mínim}: aquí $x_V=\dfrac{8}{2}=4$, amb
        $C(4)=16-32+20=4$. Interessa perquè indica amb quants recursos (monitors,
        unitats…) el cost és el més baix possible: és el \emph{punt òptim} de decisió.
\end{sollist}
\end{solucions}
